% Created 2025-12-17 mié 14:28
% Intended LaTeX compiler: pdflatex
\documentclass[11pt]{article}
\usepackage[utf8]{inputenc}
\usepackage[T1]{fontenc}
\usepackage{graphicx}
\usepackage{longtable}
\usepackage{wrapfig}
\usepackage{rotating}
\usepackage[normalem]{ulem}
\usepackage{amsmath}
\usepackage{amssymb}
\usepackage{capt-of}
\usepackage{hyperref}
\usepackage{minted}

\usepackage{parskip}
\usepackage[utf8]{inputenc} %% For unicode chars
\usepackage{placeins}
\usepackage[margin=2.50cm]{geometry}
\usepackage[T1]{fontenc}
\usepackage{mathpazo}
\linespread{1.05}
\usepackage[scaled]{helvet}
\usepackage{courier}
\hypersetup{colorlinks=true,linkcolor=blue}
\RequirePackage{fancyvrb}
\DefineVerbatimEnvironment{verbatim}{Verbatim}{fontsize=\small,formatcom = {\color[rgb]{0.5,0,0}}}
\usepackage{mdframed}
\BeforeBeginEnvironment{minted}{\begin{mdframed}}
\AfterEndEnvironment{minted}{\end{mdframed}}
\setcounter{secnumdepth}{3}
\date{}
\title{}
\hypersetup{
 pdfauthor={},
 pdftitle={},
 pdfkeywords={},
 pdfsubject={},
 pdfcreator={Emacs 27.1 (Org mode 9.6.18)}, 
 pdflang={Spanish}}
\begin{document}

\setcounter{tocdepth}{3}
\tableofcontents

\setlength\parindent{10pt}


\section{To-Do App}
\label{sec:org99cd46c}
Database design for a To-Do App with user registration and role
management, refined version of your database schema, including a table for
categories and foreign key constraints for relational integrity:

\subsubsection{Refined Database Schema}
\label{sec:orga4bb73c}
\begin{itemize}
\item Table: \texttt{todos}
\label{sec:org0c4d9dd}
\begin{minted}[]{sql}
CREATE TABLE todos (
  id INT PRIMARY KEY AUTO_INCREMENT,
  user_id INT NOT NULL,
  title VARCHAR(255) NOT NULL,
  description TEXT,
  category_id INT,
  done BOOLEAN NOT NULL DEFAULT FALSE,
  created_at TIMESTAMP NOT NULL DEFAULT CURRENT_TIMESTAMP,
  updated_at TIMESTAMP,
  FOREIGN KEY (user_id) REFERENCES users(id),
  FOREIGN KEY (category_id) REFERENCES categories(id)
);
\end{minted}

\item Table: \texttt{users}
\label{sec:org3b9da34}
\begin{minted}[]{sql}
CREATE TABLE users (
  id INT PRIMARY KEY AUTO_INCREMENT,
  username VARCHAR(255) NOT NULL UNIQUE,
  password VARCHAR(255) NOT NULL,
  role VARCHAR(255) NOT NULL DEFAULT 'user'
);
\end{minted}

\item Table: \texttt{categories}
\label{sec:orgc53f2bb}
\begin{minted}[]{sql}
CREATE TABLE categories (
  id INT PRIMARY KEY AUTO_INCREMENT,
  name VARCHAR(255) NOT NULL UNIQUE
);
\end{minted}
\end{itemize}

\subsubsection{Additional Considerations}
\label{sec:orgac8fe7b}
\begin{enumerate}
\item \textbf{Categories Table}: A separate table for categories allows more
flexibility and normalization. Users can create custom categories, or
you can pre-populate it with common categories.

\item \textbf{Foreign Keys}: These ensure relational integrity. For example,
\texttt{user\_id} in the \texttt{todos} table references \texttt{id} in the \texttt{users} table.

\item \textbf{Timestamps}: Consider automatically updating the \texttt{updated\_at}
timestamp. In MySQL, you can use \texttt{ON UPDATE CURRENT\_TIMESTAMP}.

\item \textbf{Password Storage}: Ensure you store hashed passwords, not plain
text. Use PHP's \texttt{password\_hash()} and \texttt{password\_verify()} for
handling passwords.

\item \textbf{Indexes}: Depending on your queries, consider adding indexes to
frequently searched fields for performance optimization.

\item \textbf{Roles Management}: If your application's roles have different
permissions or capabilities, consider having a separate \texttt{roles} table
and link it to the \texttt{users} table. This allows for easier management
and scalability of roles.
\end{enumerate}

\subsubsection{PHP OOP Approach}
\label{sec:org1345282}
For the PHP backend, consider the following OOP structure:

\begin{itemize}
\item \textbf{Classes for Database Tables}: Create classes like \texttt{Todo}, \texttt{User},
\texttt{Category} each encapsulating CRUD operations for their respective
tables.
\item \textbf{Session Management}: Use PHP sessions for handling user logins.
\item \textbf{Authentication and Authorization}: Implement methods in your \texttt{User}
class for handling user authentication and role-based authorization.
\item \textbf{Validation and Sanitization}: Include methods for validating and
sanitizing input data to ensure security and data integrity.
\item \textbf{Error Handling}: Robust error handling mechanisms to manage and
respond to any exceptions or errors that occur during database
operations.
\end{itemize}

\subsubsection{Example: PHP Class for \texttt{Todo}}
\label{sec:org97d89a6}
Here's a basic structure for the \texttt{Todo} class:

\begin{minted}[]{php}
class Todo {
    private $pdo;

    public function __construct($pdo) {
        $this->pdo = $pdo;
    }

    public function create($userId, $title, $description, $categoryId) {
        // Insert a new todo into the database
    }

    public function markDone($todoId) {
        // Mark a specific todo as done
    }

    // Additional CRUD methods...
}
\end{minted}

Each class (\texttt{Todo}, \texttt{User}, \texttt{Category}) will follow a similar pattern,
with methods for each CRUD operation.

\subsubsection{Web Interface}
\label{sec:org1b1f087}
For the web interface, you'll create HTML forms and PHP scripts to
interact with these classes. For example, forms for adding a todo, user
registration, login, etc., and corresponding PHP scripts to process
these forms and call the appropriate class methods.

\subsubsection{Security}
\label{sec:org215c65e}
\begin{itemize}
\item Use prepared statements to prevent SQL injection.
\item Implement CSRF tokens in forms to prevent cross-site request forgery.
\item Hash passwords using \texttt{password\_hash()} for storage.
\item Implement proper session management for user authentication.
\end{itemize}

By combining these database structures and OOP principles in PHP, you'll
create a robust, scalable, and secure To-Do application.
\end{document}
